\chapter{The Fourier--Deligne Transform}
\label{chap:fourier}

Fourier transform replaces a cohomological trace by a stalk in a family.
We normalize by the shift \([1]\).  Consequently the trace function has
a minus sign, the inverse-character composite has no reflection and
equals a Tate twist, the same-character square has a reflection, and
Plancherel has the factor \(q^n\) over \(\mathbf F_{q^n}\).

\section{The Artin--Schreier kernel}

Fix a nontrivial additive character
\(\psi:k_0\to\Qbarl^\times\).  For the finite \'etale Galois cover
\[
  z\longmapsto z^q-z
\]
of \(\A^1_{k_0}\), let \(\cL_\psi\) be the rank-one isotypic summand
corresponding to \(\psi\).

\begin{proposition}[Trace of the Artin--Schreier sheaf]
\label{prop:as-trace}
The idempotent may be normalized so that, for
\(x\in\mathbf F_{q^n}\),
\[
  t_{\cL_\psi}^{(n)}(x)
  =\psi\!\left(\Tr_{\mathbf F_{q^n}/\Fq}(x)\right).
\]
Moreover
\[
  m^*\cL_\psi\simeq p_1^*\cL_\psi\otimes p_2^*\cL_\psi,
  \qquad
  \cL_\psi^\vee\simeq\cL_{\psi^{-1}}.
\]
\end{proposition}

\begin{proof}
If \(z^q-z=x\), then
\[
 z^{q^n}-z=x+x^q+\cdots+x^{q^{n-1}}
 =\Tr_{\mathbf F_{q^n}/\Fq}(x).
\]
Thus Frobenius acts on the torsor fiber by the displayed translation.
The chosen isotypic projector gives the asserted scalar.  Additivity of
the trace gives the tensor identity, and inversion of the character
gives the dual.
\end{proof}

\begin{convention}
\label{conv:as-sign}
The trace identity in \cref{prop:as-trace} is our sign test.  Replacing
\(\psi\) by \(\psi^{-1}\) changes where a reflection appears in Fourier
inversion but changes no weight estimate.
\end{convention}

For \(a\in\mathbf F_{q^n}\), write
\(\cL_{\psi_a}=m_a^*\cL_\psi\), where \(m_a(x)=ax\).

\section{Definition, stalks, and trace functions}

On \(A_x\times A_y\), with projections \(p_x,p_y\) and
\(\mu(x,y)=xy\), define
\[
 \FT_\psi(\cK_0)
 =
 Rp_{y!}\bigl(p_x^*\cK_0\otimes\mu^*\cL_\psi\bigr)[1].
\tag{10.1}\label{eq:fourier-definition}
\]

\begin{proposition}[Geometric stalk formula]
\label{prop:fourier-stalk}
For a geometric point \(\bar a\to A_y\),
\[
  \FT_\psi(\cK)_{\bar a}
  \simeq
  R\Gamma_c(A_x,\cK\otimes\cL_{\psi_a})[1],
\]
functorially and compatibly with Frobenius.
\end{proposition}

\begin{proof}
This is compact-support base change for \(p_y\); the kernel restricts on
the fiber to \(\cL_{\psi_a}\).
\end{proof}

\begin{corollary}
\label{cor:fourier-cohomology-sheaves}
If \(\cK\) is an ordinary sheaf, then
\[
 \cH^i(\FT_\psi\cK)_{\bar a}
 \simeq H_c^{i+1}(A_x,\cK\otimes\cL_{\psi_a}).
\]
In particular, concentration of the groups on the right in degree \(1\)
makes the transform an ordinary sheaf in degree \(0\).
\end{corollary}

For a Weil complex put
\[
 t_{\cK}^{(n)}(x)
 =\sum_i(-1)^i
 \Tr(F_x\mid\cH^i(\cK)_{\bar x}).
\]

\begin{proposition}[Fourier transform of the trace function]
\label{prop:fourier-trace-function}
For \(y\in\mathbf F_{q^n}\),
\[
 t_{\FT_\psi(\cK)}^{(n)}(y)
 =
 -\sum_{x\in\mathbf F_{q^n}}
 t_{\cK}^{(n)}(x)
 \psi\!\left(\Tr_{\mathbf F_{q^n}/\Fq}(xy)\right).
\tag{10.2}\label{eq:finite-fourier-trace}
\]
\end{proposition}

\begin{proof}
Apply the trace formula to the stalk complex in
\cref{prop:fourier-stalk}.  The shift \([1]\) contributes the minus
sign.
\end{proof}

\section{Fourier inversion}

Let \(a(x)=-x\).

\begin{theorem}[Fourier inversion]
\label{thm:fourier-inversion}
There are functorial Frobenius-equivariant isomorphisms
\[
 \FT_{\psi^{-1}}\FT_\psi(\cK_0)\simeq\cK_0(-1),
\tag{10.3}\label{eq:fourier-inversion}
\]
\[
 \FT_\psi^2(\cK_0)\simeq a^*\cK_0(-1).
\]
\end{theorem}

\begin{proof}
The inverse-character composite has kernel
\[
 Rp_{xz!}\cL_\psi(y(x-z))[2].
\]
For \(x\ne z\), all compactly supported cohomology of the nontrivial
Artin--Schreier sheaf on \(A_y\) vanishes.  For \(x=z\),
\[
 R\Gamma_c(A_y,\Qbarl)=\Qbarl(-1)[-2].
\]
Thus the kernel is the extension by zero of \(\Qbarl(-1)\) from the
diagonal.  With \(\psi\) in both transforms the phase is \(y(x+z)\), so
the kernel is supported on the antidiagonal.  On functions:
\[
 \widehat{\widehat f_\psi}_{\psi^{-1}}(z)=q^nf(z),
 \qquad
 \widehat{\widehat f_\psi}_{\psi}(z)=q^nf(-z).
\]
\end{proof}

\begin{corollary}
\label{cor:fourier-equivalence}
\(\FT_\psi\) is an equivalence with quasi-inverse
\(\FT_{\psi^{-1}}(1)\), hence is conservative.
\end{corollary}

\section{Constructibility and perversity}

\begin{foundation}
\label{found:fourier-perverse}
\textbf{Perverse exactness capsule.}
On the affine line,
\[
 \FT_\psi:\operatorname{Perv}(A_x,\Qbarl)
 \xrightarrow{\sim}\operatorname{Perv}(A_y,\Qbarl)
\]
is exact.  Stalk bounds follow from affine Artin vanishing and
\cref{prop:fourier-stalk}; costalk bounds follow by Verdier duality and
\[
 \mathbb D\FT_\psi(\cK)
 \simeq\FT_{\psi^{-1}}(\mathbb D\cK)(1).
\]
Fourier inversion gives essential surjectivity.  See
\cite{LaumonFourier},
\cite[Chapter~I, Section~5]{KiehlWeissauer}, and
\cite{StanfordL23}.
\end{foundation}

\section{Plancherel}

For \(f:\mathbf F_{q^n}\to\mathbf C\), set
\(\|f\|_{2,n}^2=\sum_x|f(x)|^2\).

\begin{theorem}[Plancherel]
\label{thm:fourier-plancherel}
For every constructible Weil complex,
\[
 \|\iota t_{\FT_\psi(\cK)}^{(n)}\|_{2,n}^2
 =
 q^n\|\iota t_{\cK}^{(n)}\|_{2,n}^2.
\tag{10.4}\label{eq:plancherel}
\]
\end{theorem}

\begin{proof}
Expand the square in \cref{eq:finite-fourier-trace}.  Character
orthogonality gives
\[
 \sum_y\psi_n((x-x')y)
 =\begin{cases}q^n,&x=x',\\0,&x\ne x'.\end{cases}
\]
\end{proof}

\section{\texorpdfstring{\(L^2\)}{L2} growth and the largest weight}

The next result is the determinant-weight radius theorem.  Its delicate
point is not Pringsheim's theorem by itself: positivity does not by
itself prevent cancellation in a rational function.  The required
noncancellation is an isotypic monodromy calculation.

\begin{foundation}
\begin{theorem}[\(L^2\)-growth criterion]
\label{thm:l2-weight-criterion}
Let \(C_0\) be a geometrically connected smooth curve and let
\(\cG_0\ne0\) be a constructible \(\iota\)-mixed Weil sheaf.  Assume
that no maximal-weight subquotient is punctual and that
\(H_c^0(C,\cG)=0\).  If \(w_{\max}\) is the largest constituent weight,
then
\[
 \limsup_{n\to\infty}\frac1n
 \log_q\|\iota t_{\cG}^{(n)}\|_{2,n}^2
 =w_{\max}+1.
\tag{10.5}\label{eq:l2-weight-growth}
\]
The same formula holds after restricting \(n\) to multiples of any
fixed positive integer.
\end{theorem}

\smallskip
\textbf{Determinant-weight proof chain.}
Semisimplification does not change trace functions.  Write its
nonpunctual part as a direct sum of pure irreducible constituents
\(\cM_\lambda\), and form the conjugate sheaf constituentwise:
\[
 \cM_\lambda^c
 =\cM_\lambda^\vee\otimes\cL_{b_\lambda},
 \qquad
 \iota(b_\lambda)=q^{w_\lambda}.
\]
Then the local eigenvalue multiset of \(\cG^c\) is the complex conjugate
of that of \(\cG\).  If
\[
 S_n=\|\iota t_{\cG}^{(n)}\|_{2,n}^2,
\]
the trace formula gives
\[
 \exp\!\left(\sum_{n\ge1}S_n\frac{T^n}{n}\right)
 =L(C_0,\cG_0\otimes\cG_0^c,T).
\tag{10.6}\label{eq:l2-L-function}
\]

Pass to a finite base extension for which
\cref{prop:positive-degree-center} supplies a central degree-one element
\(z\), and decompose the maximal-weight part into its \(z\)-isotypic
blocks.  On a block \(M\), \(z\) acts by a scalar.  In
\[
 M\otimes M^c\simeq\End(M)\otimes L_{b_M},
\]
the geometrically constant quotient is
\[
 \End_{\pi_1(C)}(M)\otimes L_{b_M}.
\]
Its multiplicity
\[
 m_M=\dim\End_{\pi_1(C)}(M)
\]
is strictly positive.  Grouping geometrically isomorphic blocks, the
diagonal multiplicity is the positive Gram number
\[
 \sum_{\Gamma}
 \dim\End_{\pi_1(C)}
 \left(\bigoplus_{\lambda\in\Gamma}M_\lambda\right)>0.
\tag{10.7}\label{eq:positive-diagonal-multiplicity}
\]
Top-degree duality adds \((-1)\).  The diagonal quotient therefore has
cohomological weight \(2w_{\max}+2\), equivalently its Frobenius
eigenvalues have absolute-value exponent \(w_{\max}+1\).

The determinant-weight lemma of Kiehl--Weissauer compares, one
\(z\)-isotypic block at a time, the divisor of the cohomological
determinant in \eqref{eq:l2-L-function}.  Its proof uses the central
decomposition and the exact invariant--coinvariant sequence:
off-diagonal blocks pair with their conjugates, whereas the order in
the diagonal block is the positive integer in
\eqref{eq:positive-diagonal-multiplicity}.  Thus this top-degree factor
survives in the reduced determinant quotient.  The hypotheses exclude
the two extra terms in that exact sequence: a maximal punctual block and
a compactly-supported degree-zero block.  Lower-weight isotypes are
holomorphic on the outermost circle.  Consequently the reduced
logarithmic derivative has radius
\[
 q^{-(w_{\max}+1)}.
\]
This is the noncancellation step; coefficient positivity is not used as
a substitute for it.

Cauchy--Hadamard gives
\[
 \limsup_n S_n^{1/n}=q^{w_{\max}+1}.
\]
Replacing Frobenius by a fixed power proves descent from the finite base
extension and the assertion about multiples.  See
\cite[Chapter~I, Sections~2.12--3.2]{KiehlWeissauer} and
\cite{StanfordL24}.  This result depends on rank-one purity and the
semisimple finite-center monodromy lemma, but not on the curve estimate
or the global direct-image theorem.
\end{foundation}

\begin{warning}
\label{warn:l2-needs-mixed}
\Cref{eq:l2-weight-growth} is false as a general definition of weight.
A punctual constituent lacks the extra \(q^n\) supplied by the number
of points, and without mixedness there is no maximal pure central
isotype.
\end{warning}

\begin{formalization}
\label{form:fourier-interface}
\textbf{Fourier interface.}
\begin{enumerate}
\item Artin--Schreier kernel with its local trace normalization.
\item Fourier transform, stalk base change, and the shifted trace
formula.
\item The inversion identities
\[
 \FT_{\psi^{-1}}\FT_\psi\simeq(-1),
 \qquad
 \FT_\psi^2\simeq a^*(-1).
\]
\item Constructibility, perverse exactness, and Plancherel.
\item The \(L^2\)-growth theorem as a separate determinant-weight
capability, with its explicit nonpunctual and \(H_c^0\) hypotheses.
\end{enumerate}
\end{formalization}
